\newcommand{\DoNotLoadEpstopdf}{}
\pdfinfoomitdate=1
\pdftrailerid{}
\pdfsuppressptexinfo=-1
\documentclass[11pt,letterpaper]{article}
\usepackage[T1]{fontenc}
\usepackage{lmodern}
\usepackage{amsmath,amssymb,amsthm,mathtools}
\usepackage[margin=1in]{geometry}
\usepackage{booktabs,array}
\usepackage{microtype}
\usepackage[hidelinks]{hyperref}
\hypersetup{pdftitle={Report 163 Growth bounds for maximal families of disjoint partitions},pdfauthor={},pdfsubject={Rigorous growth bounds and correction of asymptotic fits for OEIS A068598}}
\setlength{\parskip}{4pt}
\setlength{\parindent}{1.2em}
\setlength{\emergencystretch}{2em}
\setcounter{tocdepth}{1}
\numberwithin{equation}{section}
\newtheorem{theorem}{Theorem}[section]
\newtheorem{lemma}[theorem]{Lemma}
\newtheorem{proposition}[theorem]{Proposition}
\newtheorem{corollary}[theorem]{Corollary}
\theoremstyle{remark}\newtheorem{remark}[theorem]{Remark}
\newcommand{\floor}[1]{\left\lfloor#1\right\rfloor}
\newcommand{\ceil}[1]{\left\lceil#1\right\rceil}
\title{Growth bounds for maximal families of disjoint partitions\\\large A correction of asymptotic fits for OEIS A068598}
\author{Report 163}
\date{3 October 2026}
\begin{document}
\maketitle
\begin{abstract}
Let $a(n)$ count inclusion-maximal unordered families of strict partitions
of $n$ in which no part occurs twice anywhere in the family. We prove
\[
 \frac15\le\liminf_{n\to\infty}\frac{\log a(n)}{n\log n}
 \le\limsup_{n\to\infty}\frac{\log a(n)}{n\log n}\le\frac14.
\]
The upper bound is elementary and explicit:
$a(n)\le 3^{n-1}n^{\lfloor(n-2)/4\rfloor}$ for $n\ge2$.
It follows from a recoverable-maximum encoding and the distinct-parts
mass inequality credited to Joel Spencer by Nathanson. The lower bound
injects ordinary queens configurations into maximal families and uses
the published lower bound of Luria and Simkin. A separate elementary
construction gives the weaker lower constant $1/7$.
Both positive-parameter models $A\exp(Bn^c)$ with $c>1$ listed in the
OEIS entry have ratio $a(n)/(A\exp(Bn^c))\to0$.
The first-passage inverse has leading-scale constants between $4$ and $5$,
without assuming monotonicity of $a(n)$. Neither existence of a limiting
coefficient, a multiplicative asymptotic equivalent, nor a transseries
is established. A bounded exact companion checks conventions and the
finite encodings independently of the asymptotic argument.
\end{abstract}
\tableofcontents

\section{The objects and the main conclusion}
For an integer $n\ge1$, let $H_n$ be the finite hypergraph with vertex set
$[n]=\{1,\ldots,n\}$ and edge set
\[
 \mathcal P_n=\left\{B\subseteq[n]:\sum_{b\in B}b=n\right\}.
\]
An edge is a strict integer partition, represented as a set of positive
parts. A family of pairwise disjoint edges is a matching. It is
\emph{inclusion-maximal} if no further edge can be added disjointly.
The edges and the family itself are unordered. Edge sizes may differ
within one family. Let $a(n)$ denote the number of these maximal
matchings and set $a(0)=1$ as the separate sequence convention.
Throughout, $\log$ is the natural logarithm.

This is the interpretation of OEIS A068598 supported by its title,
examples, and maximal-clique Python code \cite{oeis}. The graph comment
should say \emph{maximal cliques}: counting all cliques would give a
different sequence. Maximality does not mean maximum cardinality.
For instance the three families at $n=8$ are
\[
\begin{split}
 &\{\{8\},\{1,7\},\{2,6\},\{3,5\}\},\\
 &\{\{8\},\{1,2,5\}\},\\
 &\{\{8\},\{1,3,4\},\{2,6\}\}.
\end{split}
\]
Their cardinalities differ. There is no quotient by a symmetry of labels.

\begin{theorem}\label{thm:main}
For every integer $n\ge2$,
\begin{equation}\label{eq:upper}
 a(n)\le 3^{n-1}n^{\lfloor(n-2)/4\rfloor}.
\end{equation}
If $Q(k)$ counts ordinary nonattacking placements of $k$ queens on a
labeled $k\times k$ board and $k=\lfloor(n-1)/5\rfloor\ge1$, then
\begin{equation}\label{eq:lowerfinite}
 a(n)\ge Q(k).
\end{equation}
Consequently,
\begin{equation}\label{eq:logbounds}
 \frac15\le\liminf_{n\to\infty}\frac{\log a(n)}{n\log n}
 \le\limsup_{n\to\infty}\frac{\log a(n)}{n\log n}\le\frac14.
\end{equation}
In particular $\log a(n)=\Theta(n\log n)$. More precisely, the lower
bound obtained from the queens theorem is the one-sided estimate
\begin{equation}\label{eq:linear}
 \log a(n)\ge \frac n5\log n-\frac{\log5+3}{5}n+o(n).
\end{equation}
\end{theorem}
The meaning of \eqref{eq:linear} is that some error sequence $r_n=o(n)$
makes the displayed inequality true for all sufficiently large $n$.
It is not an equality or a determination of the linear term of
$\log a(n)$. The upper bound is self-contained; the stronger lower
constant uses precisely the external theorem stated in Section~\ref{sec:queens}.

\section{An encoding with at most one quarter free labels}
\label{sec:upper}
Every maximal matching contains $\{n\}$. Indeed, no other edge of $H_n$
contains $n$, since every part is positive. Thus $\{n\}$ could be added
to any matching lacking it. Remove this forced edge and write the
remaining matching as $\mathcal M$. Let $m=|\mathcal M|$ and let $v$
be the number of labels used by its edges. Every remaining edge has at
least two parts, so $v\ge2m$.

For each $B\in\mathcal M$, retain its smallest part $a_B$ as an
\emph{anchor}, omit its largest part $z_B$, and call its other parts
\emph{free}. Define the disjoint subsets of $[n-1]$
\[
 A=\{a_B:B\in\mathcal M\},\qquad
 F=\bigcup_{B\in\mathcal M}(B\setminus\{a_B,z_B\}),
\]
and define $f:F\to A$ by assigning each free part to the anchor of its
edge. Put $q=|F|=v-2m$.

\begin{lemma}\label{lem:code}
The code $(A,F,f)$ determines $\mathcal M$ uniquely.
\end{lemma}
\begin{proof}
For each $a\in A$, recover the corresponding edge as
\begin{equation}\label{eq:decode}
 \{a\}\ \cup\ f^{-1}(a)\ \cup\
 \left\{n-a-\sum_{x\in f^{-1}(a)}x\right\}.
\end{equation}
For a code that came from a matching, the final part is its original
largest part. This also recovers two-part edges, whose fibers are empty.
The empty matching has the unique empty code. Unused labels need no
additional information. Arbitrary codes may fail to reconstruct valid
edges or disjoint families; the count below only overcounts by including
them.
\end{proof}

The used labels are $v$ distinct positive integers and their sum is
$mn$. Therefore
\begin{equation}\label{eq:mass}
 mn\ge1+2+\cdots+v=\frac{v(v+1)}2.
\end{equation}
Nathanson \cite[p.~108]{nathanson} explicitly attributes the fixed-arity
version of this summation argument to Joel Spencer. We use the same
ingredient for variable edge sizes and combine it with the encoding;
we make no claim to have originated the mass inequality.

Using \eqref{eq:mass} gives
\begin{align}
 q=v-2m
 &\le v-\frac{v(v+1)}n\notag\\
 &=\frac{(n-1)^2}{4n}
   -\frac1n\left(v-\frac{n-1}{2}\right)^2
 \le\frac{(n-1)^2}{4n}.\label{eq:qbound}
\end{align}
For integers $n\ge2$,
\begin{equation}\label{eq:floor}
 \floor{\frac{(n-1)^2}{4n}}=\floor{\frac{n-2}{4}}.
\end{equation}
Indeed, the left argument equals $(n-2)/4+1/(4n)$; the fractional
part of $(n-2)/4$ belongs to $\{0,1/4,1/2,3/4\}$ and adding
$1/(4n)\le1/8$ does not cross the next integer.

There are $3^{n-1}$ choices of a disjoint pair $(A,F)$, assigning each
label to $A$, $F$, or neither. For a pair supporting an actual code,
the number of maps is at most $m^q\le n^q$, and
\eqref{eq:qbound}--\eqref{eq:floor} bound $q$. When $A=F=\varnothing$
there is exactly one empty function; when $A=\varnothing$ and
$F\ne\varnothing$ there is none. Thus no convention about $0^0$ is
needed. Lemma~\ref{lem:code} now proves \eqref{eq:upper}, and hence
\begin{equation}\label{eq:logupper}
 \log a(n)\le\tfrac14n\log n+(n-1)\log3.
\end{equation}
Only the initial forced-singleton step used maximality. In particular,
the code is also injective on every matching without the singleton.

\section{Recoverable bases from queens configurations}\label{sec:queens}
A queens configuration on a labeled $k\times k$ board is represented by
a permutation $\pi$ of $[k]$ such that both $i+\pi(i)$ and $i-\pi(i)$
are pairwise distinct. These are ordinary diagonals, with no reduction
modulo $k$, and configurations are not identified under board symmetries.
The sole external asymptotic input is the following published theorem.

\begin{theorem}[Luria and Simkin]\label{thm:queens}
As $k\to\infty$ through all positive integers,
\begin{equation}\label{eq:queens}
 Q(k)\ge \bigl((1-o(1))k/e^3\bigr)^k.
\end{equation}
\end{theorem}
The primary version used is Theorem~1.1 on page~1 of
arXiv:2105.11431v2 \cite{queens}. Its page~2 also explicitly gives
$\log Q(k)/(k\log k)\to1$, using $Q(k)\le k!$ for the other direction.
The theorem is for ordinary queens through all integers. The use of
toroidal placements inside its proof imposes no coprimality condition
on the theorem used here. That external proof is not reproduced in
this report.

Fix $n$ and $k=\lfloor(n-1)/5\rfloor\ge1$. Then $n>5k$. Given a
queens permutation $\pi$, define
\begin{equation}\label{eq:triples}
 B_i(\pi)=\{i,\ k+\pi(i),\ n-k-i-\pi(i)\},
 \qquad 1\le i\le k.
\end{equation}
Each triple sums to $n$. Its three value classes lie in
\[
 [1,k],\qquad[k+1,2k],\qquad[n-3k,n-k-2],
\]
respectively, and $n-3k>2k$. Thus the classes are disjoint and each
triple is strict. The first two classes have no repetitions because
$\pi$ is a permutation; the third has no repetitions because the sums
$i+\pi(i)$ are distinct. The triples therefore form a matching.
Only one diagonal condition is needed for this embedding; ordinary
queens provide a subset for which the required count is known.

\begin{lemma}\label{lem:extension}
Choose any maximal extension in $H_n$ of each matching
$\{B_i(\pi):1\le i\le k\}$. Distinct queens permutations receive
distinct extensions, regardless of these choices.
\end{lemma}
\begin{proof}
A maximal extension exists because the edge set of $H_n$ is finite:
add a disjoint edge whenever possible until none remains. Extensions
can include edges of all sizes. If a matching contained the bases for
both $\pi$ and $\sigma$, then its unique edge containing each
$i\in[1,k]$ would equal both $B_i(\pi)$ and $B_i(\sigma)$. Each such
edge has exactly one member of $[k+1,2k]$, so $\pi(i)=\sigma(i)$ for
every $i$. This proves injectivity.
\end{proof}
Lemma~\ref{lem:extension} proves \eqref{eq:lowerfinite}.
The cases $k=2,3$, where $Q(k)=0$, are harmless vacuous inequalities.
For large $k$, taking logarithms in \eqref{eq:queens} yields
\[
 \log a(n)\ge k\log k-3k+o(k).
\]
Writing $n=5k+r$ with $r\in\{1,2,3,4,5\}$ gives
\[
 k\log k=\frac n5\log n-\frac{\log5}{5}n+O(\log n).
\]
It follows that \eqref{eq:linear} holds for every sufficiently large
$n$, rather than only a subsequence. Together with
\eqref{eq:logupper}, this completes Theorem~\ref{thm:main}.

\section{A self contained elementary lower bound}
The $1/5$ constant depends on Theorem~\ref{thm:queens}, but an
elementary construction already proves factorial-logarithmic growth.

\begin{proposition}\label{prop:greedy}
If positive integers $k,m$ satisfy $m\ge2k$ and $3k+2m<n$, then
\begin{equation}\label{eq:greedyfinite}
 a(n)\ge\prod_{i=1}^k(m-2(i-1))\ge(m-2k+2)^k.
\end{equation}
In particular, for $n\ge8$ and $k=\lfloor(n-1)/7\rfloor$,
\begin{equation}\label{eq:factorial}
 a(n)\ge 2^k k!.
\end{equation}
Consequently,
\begin{equation}\label{eq:seventh}
 \liminf_{n\to\infty}\frac{\log a(n)}{n\log n}\ge\frac17
\end{equation}
without invoking the queens theorem.
\end{proposition}
\begin{proof}
Let $A=[1,k]$ and $B=[k+1,k+m]$. Choose $b_i\in B$ successively,
requiring all $b_i$ and all sums $i+b_i$ to be distinct. At step $i$,
repeating an earlier $b_j$ excludes at most $i-1$ choices. An equality
$i+b_i=j+b_j$ excludes at most another $i-1$ choices. At least
$m-2(i-1)$ possibilities remain, proving the product count.
Set $c_i=n-i-b_i$. Then
\[
 c_i\ge n-2k-m>k+m.
\]
The triples $\{i,b_i,c_i\}$ are strict partitions of $n$ with disjoint
value ranges. Within those ranges there are no repeated values, the
last because the sums $i+b_i$ are distinct. Each base is a matching,
and every anchor in $A$ is used. The argument of
Lemma~\ref{lem:extension}, with middle range $B$, proves that distinct
choices cannot share a maximal extension. This proves
\eqref{eq:greedyfinite}.

Set $m=2k$ and retain the full product rather than its smallest
factor. If $n>7k$, then
\[
 \prod_{i=1}^k(2k-2(i-1))=2^k k!,
\]
which proves \eqref{eq:factorial} for $k=\lfloor(n-1)/7\rfloor$.
The elementary estimate
$\log(k!)=k\log k-k+O(\log k)$ gives
\[
 \log a(n)\ge\frac n7\log n+O(n),
\]
and proves \eqref{eq:seventh}.
\end{proof}
Together with \eqref{eq:upper}, this gives a wholly elementary proof
of $\log a(n)=\Theta(n\log n)$ and of the rejection of the fits below.
The factorial bound also provides finite first-passage bounds.

\section{Why the posted asymptotic fits fail}
The inspected OEIS entry records a 2022 conjectural model with
parameters approximately $(A,B,c)=(0.57,0.062,1.54)$, and a February
2026 revision to $(0.26,0.129,1.36)$ \cite{oeis}. The following statement
rules out either model as an asymptotic equivalent, regardless of
its finite-range numerical quality.

\begin{corollary}\label{cor:fits}
For every fixed $A,B>0$ and $c>1$,
\begin{equation}\label{eq:fits}
 \frac{a(n)}{A\exp(Bn^c)}\longrightarrow0.
\end{equation}
\end{corollary}
\begin{proof}
By \eqref{eq:logupper},
\[
 \log\frac{a(n)}{A\exp(Bn^c)}
 \le\tfrac14n\log n+O(n)-Bn^c-\log A\longrightarrow-\infty,
\]
since $n\log n=o(n^c)$ for $c>1$. Exponentiation proves the claim.
\end{proof}
Only the elementary upper bound is needed for this conclusion.
It does not imply that the two expressions cannot approximate a
short data range. No regression in the companion is used to prove or
choose a growth coefficient.

\section{First passage inversion without monotonicity}
For a real threshold $y>1$ define
\[
 N(y)=\min\{n\ge0:a(n)\ge y\}.
\]
The all-integer lower bound implies that this set is nonempty for all
$y>1$. Since each $a(n)$ is finite, $N(y)\to\infty$ as $y\to\infty$.

\begin{corollary}\label{cor:inverse}
As $y\to\infty$,
\begin{equation}\label{eq:inverse}
 4\le\liminf N(y)\frac{\log\log y}{\log y}
 \le\limsup N(y)\frac{\log\log y}{\log y}\le5.
\end{equation}
In particular $N(y)=\Theta(\log y/\log\log y)$.
\end{corollary}
\begin{proof}
Write $L=\log y$. Fix $0<u<4$ and $w>5$, and choose $\varepsilon>0$
so that $(1/4+\varepsilon)u<1$ and $(1/5-\varepsilon)w>1$.
Theorem~\ref{thm:main} supplies a fixed $n_0$ such that for all
integers $n\ge n_0$,
\begin{equation}\label{eq:epsilon}
 (1/5-\varepsilon)n\log n\le\log a(n)
 \le(1/4+\varepsilon)n\log n.
\end{equation}
Let $n_- =\lfloor uL/\log L\rfloor$. For every
$n_0\le n\le n_-$, monotonicity of the elementary function $n\log n$
and \eqref{eq:epsilon} give
\[
 \log a(n)\le(1/4+\varepsilon)n_-\log n_-
 =(1/4+\varepsilon)uL(1+o(1))<L.
\]
The finitely many $n<n_0$ also have $a(n)<y$ for large $y$.
Thus $N(y)>n_-$. Next put $n_+=\lceil wL/\log L\rceil$.
The lower bound in \eqref{eq:epsilon} gives
\[
 \log a(n_+)\ge(1/5-\varepsilon)n_+\log n_+
 =(1/5-\varepsilon)wL(1+o(1))>L,
\]
so $N(y)\le n_+$. Letting $u$ increase to $4$ and $w$ decrease to $5$
proves \eqref{eq:inverse}.
\end{proof}
No comparison of $a(n)$ and $a(n+1)$ occurs in this proof. We assert
neither monotonicity of the sequence nor an exact inverse equivalent.
The inequalities do not justify an exact ceiling formula. If only the
elementary lower bound is used, the same proof gives constants $4$
and $7$ instead.

\subsection{Finite elementary inverse bounds}
Let $W$ be the positive real Lambert function: $W(t)e^{W(t)}=t$ for
$t>0$. The elementary estimates also give explicit finite bounds.
\begin{proposition}\label{prop:finiteinverse}
For every real $y>1$, with $L=\log y$,
\begin{equation}\label{eq:finiteinverse}
 \ceil{\frac{4(L+\log3)}{W(324(L+\log3))}}
 \le N(y)\le
 7\ceil{\frac{L}{W(2L/e)}}+1.
\end{equation}
\end{proposition}
\begin{proof}
The continuous function $h(x)=x\log(81x)/4$ is strictly increasing
for $x\ge1$, and $h(1)=\log3$. Its unique root $r>1$ of
$h(r)=L+\log3$ is
\[
 r=\frac{4(L+\log3)}{W(324(L+\log3))}.
\]
By \eqref{eq:logupper}, every integer $2\le n<r$ has
$\log a(n)\le h(n)-\log3<L$. The values $n=0,1$ separately have
$a(n)=1<y$. Thus $N(y)\ge\lceil r\rceil$.

For the upper bound let $s=L/W(2L/e)$. Then $s>e/2$ and
$s\log(2s/e)=L$. The function $x\log(2x/e)$ is strictly increasing
for $x\ge e/2$. With $k=\lceil s\rceil\ge2$, the elementary integral
bound $\log(k!)\ge k\log k-k$ gives
\[
 2^k k!\ge(2k/e)^k\ge e^L=y.
\]
At $n=7k+1$, \eqref{eq:factorial} therefore yields
$a(n)\ge y$, proving the right inequality without monotonicity.
\end{proof}
These are exact-real mathematical inequalities. Uncertified floating-
point evaluations of $W$ near an integer do not certify the ceilings.
The companion instead uses integer powers and factorials to produce
finite brackets at integer thresholds: it excludes every index whose
explicit upper bound is too small, and supplies an index $7k+1$ for
which $2^k k!$ reaches the threshold. The finite elementary upper
bound has leading constant $7$; the asymptotic queens estimate, without an explicit
onset here, improves the asymptotic upper constant to $5$.

\section{Finite verification and the reproducible companion}
The companion uses exact integer arithmetic and the Python standard
library. Its production module has no network, input-file, or
file-writing API. It prints deterministic JSON and imposes fixed bounds
before enumeration. The code is standalone; it does not require any
other report or unpublished working file. See the accompanying README
for complete commands and output-safety assumptions.

Strict partitions are constructed by subset-sum dynamic programming
and checked against a separate recursive generator. A direct increasing-
edge matching traversal tests maximality against \emph{every} edge,
including earlier edges not chosen. A separate maximal-clique traversal
provides another count comparison in a smaller range. Removal of the
forced singleton does not change the maximal-family count.

The exact tests cover the following.
\begin{enumerate}
\item All values $a(n)$ for $0\le n\le24$ agree with the inspected
OEIS data. The singleton, unordered-family, and inclusion-maximal
conventions are checked explicitly.
\item Every matching without its singleton through $n=24$ is encoded
and decoded, including the empty matching and empty fibers. No two
matchings have the same code. The mass inequality, the integer free-
label bound, and the explicit upper bound are checked exactly.
\item Ordinary queens are enumerated through $k=8$. For every one of
the five choices $n=5k+r$, $1\le r\le5$, actual maximal extensions
of the triple bases are constructed. Original permutations are
recovered from the extensions and distinctness is checked.
\item The upper floor identity and the five-residue inequalities are
checked through $n=100000$. Bounded greedy-construction examples
check the choice counts and disjoint triple ranges.
\item Argument validation, workload caps, first-passage scans,
intentional fault detection, and byte-identical normal and optimized
Python outputs are exercised. No correctness check relies on an
\texttt{assert} statement that disappears under \texttt{-O}.
\end{enumerate}
The first values include
\[
\begin{array}{c|rrrrrrrrr}
n&16&17&18&19&20&21&22&23&24\\\hline
a(n)&75&115&199&312&533&888&1536&2535&4608.
\end{array}
\]
The threshold command performs an exact bounded scan; failure to reach
a threshold only proves that its first passage exceeds the scanned
range. It does not extrapolate using the asymptotic bounds. A separate integer
bracket command implements the proven finite elementary bounds; its
endpoints are not asserted to be exact first-passage indices.

These computations are finite diagnostics, not proofs of the
asymptotic statements or of novelty. Earlier independent checks
through $n=22$ and a separate direct-enumeration audit through $n=24$
served the same diagnostic role. The proof consists of the explicit
arguments above together with Theorem~\ref{thm:queens}.

The release includes complete LaTeX, a deterministic PDF builder,
fixed-allowlist ZIP creation, and SHA-256 checksums. Builds and archive
creation use new output names and refuse existing destinations and
symlink components; they assume a trusted local parent directory.
Two clean PDF builds and two archive constructions are compared
byte for byte in the reported environment. Reproduction with a
different TeX installation can legitimately change PDF bytes.

\section{Source provenance and the remaining gap}
The official A068598 snapshot used here is revision \#50, dated
17 February 2026, retrieved from the official OEIS data repository on
3 October 2026. Its definition, displayed terms, and fit comments
were also checked against the public entry. A site-wide footer is
not treated as an entry revision date. Exact snapshot digests and
the inspected source locations appear in the companion's provenance
file; full third-party papers and full entry prose are not redistributed.

Nathanson's paper \cite{nathanson} was read beyond its abstract.
It studies the largest and smallest cardinalities of maximal families
when \emph{every} partition has a fixed number of parts. That is a
different extremal problem from counting unrestricted maximal
families. Its page~108 gives the Spencer attribution used above.
Filaseta's related 1996 chapter \cite{filaseta} was located and its
publisher abstract inspected. That abstract likewise concerns a
minimum cardinality at fixed arity. The full subscription chapter
was not read, so its arguments and remarks have not been excluded
as possible prior art.

Bounded exact-ID and topic searches found no primary source claiming
the unrestricted enumeration bounds proved here. Search terms
included A068598 with ``asymptotic'' or ``growth'', maximal pairwise
disjoint partitions with ``number'' or ``asymptotic'', and disjoint
partitions with ``$n\log n$''. This is a bounded literature check,
not a priority certificate. The result should be described as an
independently verified proof candidate and application of known
ingredients, rather than as a certified first proof.

The principal mathematical gap is still
\[
 \frac15\le\liminf\frac{\log a(n)}{n\log n}
 \le\limsup\frac{\log a(n)}{n\log n}\le\frac14.
\]
The report does not prove that the lower and upper limits coincide,
identify their value, establish a multiplicative equivalent, or
supply any all-orders expansion or transseries. The linear term in
\eqref{eq:linear} belongs only to a lower bound. Improving either
counting construction or showing a matching upper and lower
coefficient remains open in this treatment.

One concrete unproved direction is to count decompositions of $[4k]$
into $k$ quadruples each summing to $8k+2$. A bound
$k^{2k-o(k)}$ would, by maximal extension and recovery of the used
labels, attain coefficient $1/4$ along $n=8k+2$. It would not by
itself settle all $n$. Requiring one part from each of four consecutive
$k$-blocks leads to permutations $\pi,\sigma,\tau$ with
$i+\pi(i)+\sigma(i)+\tau(i)=2k+2$. The elementary complementary-pair
construction gives only $k!$, far short of that proposed count.
No stronger count or resulting limit is claimed here.

\begin{thebibliography}{9}
\bibitem{oeis}
The OEIS Foundation Inc., \emph{The On-Line Encyclopedia of Integer
Sequences}, A068598, entry by Naohiro Nomoto with later contributions.
Revision \#50, 17 February 2026; inspected 3 October 2026.
\href{https://oeis.org/A068598}{oeis.org/A068598}.
Official data: \href{https://github.com/oeis/oeisdata/blob/main/seq/A068/A068598.seq}{A068598.seq}.

\bibitem{nathanson}
M.~B. Nathanson, \emph{Largest and smallest maximal sets of pairwise
disjoint partitions}, Journal of Number Theory \textbf{17} (1983),
103--112. Spencer's summation argument is on p.~108.
\href{https://www.theoryofnumbers.com/melnathanson/pdfs/nath1983-49.pdf}{Author-hosted primary PDF}.

\bibitem{queens}
Z.~Luria and M.~Simkin, \emph{A lower bound for the $n$-queens problem},
Proceedings of the 2022 Annual ACM--SIAM Symposium on Discrete
Algorithms (SODA), 2022.
\href{https://doi.org/10.1137/1.9781611977073.86}{doi:10.1137/1.9781611977073.86}.
Version inspected: \href{https://arxiv.org/abs/2105.11431v2}{arXiv:2105.11431v2}
(9 July 2021), Theorem~1.1 on PDF p.~1 and consequence on p.~2.

\bibitem{filaseta}
M.~Filaseta, \emph{The smallest maximal set of pairwise disjoint
partitions}, in D.~V. Chudnovsky, G.~V. Chudnovsky, and M.~B. Nathanson
(eds.), \emph{Number Theory: New York Seminar 1991--1995}, Springer,
1996, 103--113.
\href{https://doi.org/10.1007/978-1-4612-2418-1_8}{doi:10.1007/978-1-4612-2418-1\_8}.
Publisher metadata and abstract were inspected. The full chapter
was not read.
\end{thebibliography}
\end{document}
